\section{Conclusion}
This paper has documented an open-source evaluation platform for comparing a two-level SiC inverter with an AC battery (BI-MMC) in two analyses that share every loss equation: torque--speed efficiency maps, and drive-cycle efficiency with SOC and temperature tracking. All equations of the Python implementation, the algorithms of both analyses and their assumptions are given.

Applied to six applications with equal machine, cells and installed energy, the platform shows that the AC battery cuts the semiconductor loss energy by \SIrange{26}{80}{\percent} but increases the battery loss energy by a factor of \numrange{1.6}{4.4}, because the cells carry the switched phase current. On a cycle basis, the AC battery is therefore at best marginally better (\SI{+0.09}{\pp} for the off-road machine and \SI{+0.10}{\pp} for the ship, where nearly all energy is delivered below base speed and the installed energy is large relative to the power) and clearly worse for the passenger car and city bus ($-0.49$ and $-0.57$~pp), the long-haul truck ($-0.83$~pp) and the regional train ($-1.74$~pp). Low SOC aggravates the losses of the AC battery. The answer to the title question is thus: only in low-speed, low-C-rate duty, and then by only about \SI{0.1}{\pp}.

The platform also shows which levers could change this answer: a higher installed energy per unit power, a lower cell resistance, a smaller field-weakening current, and a modulation that reduces the cell current. The platform, the application presets, the data files and the scripts that generate every figure and table of this paper are provided in the repository, so that the study can be repeated with other machines, cells, devices and duty cycles.
